\subsection{Ordinary fourth moments for genuine singleton-root inputs}\label{proof:Q5-C087}
\paragraph{Q87.0 Singleton-root fourth moments and their outer transfer}\label{part:Q87.0}

Work on the finite probability space with a real skew kernel. Put

\[
 A=J+K,\quad T=J-K=A^*,\quad g=K1,\quad p=1+g,\quad q=1-g,
\]
\[
\begin{gathered}
f_n=A^n1,\quad \bar f_n=T^n1,\quad W_n=A \bar f_{n-1},\qquad\\
I_q(U)=E[q^{-1/3}|U|^{8/3}].
\end{gathered}
\]

The original offcenter root has precisely one residual arm of length
\(a\), after maximal actual equal-length pairing. The terminal has
precisely the two residual arms \((1,\allowbreak  b)\), \(b\ge\allowbreak 2\). Their inner
connector has length \(L\), and the independently specified connector
from that root to the receiving root has length \(M\). Equal-arm factors are restored by positivity in the common tree assembly.

Local bound. For every

\[
 a\ge3,\qquad b\ge2,\qquad L\ge1,\qquad M\ge1,
\]

the following orientation, determined by the lengths, gives a receiving-root message \(U\)
satisfies \(I_q(U)\le\allowbreak 1\).


\begingroup\small
\setlength{\tabcolsep}{3pt}
\setlength{\LTleft}{0pt}\setlength{\LTright}{\fill}
\begin{longtable}{@{}p{0.2200\linewidth}p{0.7000\linewidth}@{}}
\caption{Ordinary fourth moments for genuine singleton-root inputs: parameter values.}\label{tab:parameters:q5-c087:1}\\
\toprule
Entry & Actual input \(Y\) at the singleton root \\
\midrule\endfirsthead
\toprule
Entry & Actual input \(Y\) at the singleton root \\
\midrule\endhead
\bottomrule\endfoot
\(L=\allowbreak 1\), every \(a\ge\allowbreak 3\) & \(f_aT(p\bar f_b)\); Outer message: \(T^M Y\). \\
\addlinespace[2pt]
\(a=\allowbreak 3,\allowbreak  L=\allowbreak 2,\allowbreak  M=\allowbreak 1,\allowbreak  2\) & \(f_3TA(p\bar f_b)\); Outer message: \(T^M Y\). \\
\addlinespace[2pt]
\(a=\allowbreak 3,\allowbreak  L=\allowbreak 2,\allowbreak  M\ge\allowbreak 3\) & \(W_3T^2(qf_b)\); Outer message: \(T^M Y\). \\
\addlinespace[2pt]
\(a=\allowbreak 3,\allowbreak  L\ge\allowbreak 3\) & \(W_3T^L(qf_b)\); Outer message: \(T^M Y\). \\
\addlinespace[2pt]
\(a\ge\allowbreak 4,\allowbreak  L\ge\allowbreak 2\) & \(W_aT^L(qf_b)\); Outer message: \(T^M Y\). \\
\end{longtable}
\endgroup


For the first, second, fourth, and the \(a=\allowbreak 4\) part of the last row,
the new proof gives

\[
 EY^4\le e^{-11d/2}.
 \tag{Q87.0.1}\label{eq:Q87:0.1}
\]

For every \(a\ge\allowbreak 5,\allowbreak  b\ge\allowbreak 2,\allowbreak  L\ge\allowbreak 2\), it gives

\[
 E\{[W_aT^L(qf_b)]^4\}\le e^{-3487d/640}.
 \tag{Q87.0.2}\label{eq:Q87:0.2}
\]

Here \(d=\allowbreak Eg^2\). The new fourth-input outer bridges are

\[
\begin{gathered}
I_q(TY)^3\le e^{204d/19}(E|Y|^4)^2,\qquad\\
I_q(T^2Y)^3\le e^{87d/8}(E|Y|^4)^2.
\end{gathered}\tag{Q87.0.3}\label{eq:Q87:0.3}
\]

Consequently every short-outer row of the local bound has the common bound

\[
 I_q(U)^3\le e^{-7d/320}\qquad(M=1,2).
 \tag{Q87.0.4}\label{eq:Q87:0.4}
\]

In (Q87.0.1) rows the stronger bounds are \(e^{-5d/19}\) for \(M=\allowbreak 1\)
and \(e^{-d/8}\) for \(M=\allowbreak 2\). For \(M\ge\allowbreak 3\), the exact
Proposition~\ref{prop:ordinary-unit-inner}, Section~\ref{proof:Q5-C075}, and the unanchored transfer of Section~\ref{proof:Q5-C079} apply to the actual rows
specified above. In particular,

\[
 I_q(U)^3\le
 \begin{cases}
 (1-d)e^{-d/1024},&L=1,\\
 (1-d)e^{-3d/20},&L\ge2.
 \end{cases}
 \tag{Q87.0.5}\label{eq:Q87:0.5}
\]

For \(a=\allowbreak 3,\allowbreak  L=\allowbreak 2\), the direction is selected from the outer length \(M\).

\paragraph{Q87.1 Finite identities and inputs}\label{part:Q87.1}

Delete only zero-mass atoms. On every retained atom,
\(p(x),\allowbreak  q(x)\ge\allowbreak \mu(x)>\allowbreak 0\), because \(K(x,\allowbreak  x)=\allowbreak 0\). Every finite tree message is positive: conditional on its root, the labeling
putting all other original vertices at that same atom has positive
mass and all diagonal edge factors one. All norms, expectations, adjoints and rotation blocks use \(\mu\).
Throughout, \(EY^j\) means the moment \(E[Y^j]\), never the
power \((EY)^j\) of a mean; displayed explicit norm
squares are interpreted in the usual original Hilbert space.

Write

\[
\begin{gathered}
r=Kg,\quad s=Kr,\quad t=Ks,\quad \delta_j=\|K^j1\|_2^2,\\
\theta=13/32,\quad\kappa=19/32,\quad\tau=d-\delta_2.
\end{gathered}
\]

The finite calculus supplies

\[
\begin{gathered}
0\le d\le1/3,\quad d^2\le \delta_2\le\theta d,\quad\\
\delta_{j+1}\le\theta \delta_j,\quad\|K\|^2\le\theta,\quad\\
|Kv(x)|\le\|v\|_2,\quad \kappa d\le\tau\le2/9,
\end{gathered}\tag{Q87.1.1}\label{eq:Q87:1.1}
\]
\[
\begin{gathered}
Er=-d,\quad Es=0,\quad Et=\delta_2,\quad Egr=Ers=0,\quad\\
Egs=-\delta_2,\quad Ert=-\delta_3.
\end{gathered}\tag{Q87.1.2}\label{eq:Q87:1.2}
\]

Original row centering and antisymmetrization give, respectively,

\[
 r^2\le pqd,\qquad |r|\le pq/2.
 \tag{Q87.1.3}\label{eq:Q87:1.3}
\]

For the second inequality, use
\(r(x)=\allowbreak \frac12E_{y,\allowbreak  z}[(K(x,\allowbreak  y)-K(x,\allowbreak  z))K(y,\allowbreak  z)]\) and
\(|u-v|\le\allowbreak 1-uv\) on \([-1,\allowbreak  1]^2\). In particular \(|r|\le\allowbreak 1/2\).


Define finite polynomials and finite alternating sums by

\[
\begin{gathered}
G_n(z)=\sum_{j=0}^{n-1}z^j,\quad\\
\beta_1=\beta_2=0,\quad\\
\beta_n=\sum_{j=1}^{\lfloor(n-1)/2\rfloor}(-1)^{j-1}\delta_j.
\end{gathered}
\]

The exact finite positive convolution, for either sign, is

\[
\begin{gathered}
\sum_{j=0}^n(\pm K)^j1\\
=f_n^\pm+\sum_{j=3}^n\beta_jf_{n-j}^\pm,\qquad\\
\tau\le\beta_j\le d\quad(j\ge3),
\end{gathered}\tag{Q87.1.4}\label{eq:Q87:1.4}
\]

where \(f_n^+=\allowbreak f_n,\allowbreak  f_n^-=\allowbreak \bar f_n\). Every removed term is a nonnegative
actual path. The finite sum identity is proved in the common finite calculus.

At \(z=\allowbreak i\sqrt x\), \(0\le\allowbreak  x\le\allowbreak \theta\), write
\(G_n(z)=\allowbreak A_n+zB_n\). For \(n\ge\allowbreak 2\), both coordinates lie in
\([1-x,\allowbreak  1]\). For \(n\ge\allowbreak 4\), both are at most \(1-x+x^2\).
Also \(\|G_n(K)\|\le\allowbreak 1\) for \(n\ge\allowbreak 3\), and
\(|G_n(i\sqrt x)|^2\le\allowbreak 1-7x/15\) for \(n\ge\allowbreak 3\).
The coordinate estimates are Lemma~\ref{lem:odd-common:spectral}; the finite decay bound is (Q81.1.5) of Section~\ref{proof:Q5-C081}.
they hold at the short endpoints and for every finite integer length.

We use the following inputs.

1. Section~\ref{proof:Q5-C081}, \hyperref[part:Q81.2]{the signed-mean exponential lemma} proves, for every original real vector \(V\) and
   every \(u\ge\allowbreak 0\),

\[
   Ee^{uKV}\le
   \exp\{uE(KV)+(17/64)u^2\|V\|^2\}.
   \tag{Q87.1.5}\label{eq:Q87:1.5}
\]

   It retains the actual mean of either sign. We use it at \(u=\allowbreak 4\)
   and \(u=\allowbreak 8\), including scaled vectors. The Hermite--Bernoulli proof covers the full continuous domain.
2. Section~\ref{proof:Q5-C064}, the short-operator and all-power estimates give

\[
\begin{gathered}
\|T\|^2\le(1+32d/19)^2,\quad\\
\|T^2\|^2\le(1+d/2)^4,\quad\\
\|T^M\|\le(1-d/8)^{\lfloor M/3\rfloor}\quad(M\ge3).
\end{gathered}\tag{Q87.1.6}\label{eq:Q87:1.6}
\]

   The Section~\ref{proof:Q5-C056} bound
   \(E[\bar f_3^4/q]\le\allowbreak 1-23d/36\).
3. Section~\ref{proof:Q5-C079}, the eighth-input transfer estimates prove the eighth-input map and its row budgets:

\[
\begin{gathered}
I_q(T^M Y)^3\le\|T^M\|^2E|Y|^8R_M^2,\quad\\
R_M=E[\bar f_M^3/\sqrt q],\quad R_M^2\le1-d (M\ge3),
\end{gathered}\tag{Q87.1.7}\label{eq:Q87:1.7}
\]

   with short costs \(e^{299d/38}\), \(e^{6d}\).
4. Lemma~\ref{lem:ordinary-unit-inner-terminal} and Proposition~\ref{prop:ordinary-unit-inner} give the actual unit-inner envelopes and ordinary eighth
   for \(f_aT(p\bar f_b)\), all \(a\ge\allowbreak 3,\allowbreak  b\ge\allowbreak 2\). Section~\ref{proof:Q5-C075} gives the actual
   nonunit-inner envelopes, their paid scalar expressions, their
   covariance signs, and the ordinary eighth for
   \(W_aT^L(qf_b)\), all \(a\ge\allowbreak 3,\allowbreak  b\ge\allowbreak 2,\allowbreak  L\ge\allowbreak 2\).
5. Section~\ref{proof:Q5-C058-ordinary} gives \(E[pT^L(qf_b)]^8\le\allowbreak 1-d/12000\) for all \(b\ge\allowbreak 2,\allowbreak  L\ge\allowbreak 1\).
   Section~\ref{proof:Q5-C071} gives \(E[W_2T^L(qf_b)]^8\le\allowbreak 1-d/256\) for \(b,\allowbreak  L\ge\allowbreak 2\).


At \(d=\allowbreak 0\), \(g=\allowbreak 0\) but \(K\) need not vanish. Leaf integration
makes every actual message one. The bounds for arbitrary real inputs
also follow from their displayed norm/Jensen proofs. Every later
division by \(d\) assumes \(d>\allowbreak 0\).

\paragraph{Q87.2 Fourth-moment transfer along short outer paths}\label{part:Q87.2}

Let \(B\) be any positive original operator with row \(h=\allowbreak B1>\allowbreak 0\).
Row Cauchy, then Holder with exponents 3 and \(3/2,\allowbreak \) give

\[
\begin{gathered}
|BY|^2\le hB(|Y|^2),\qquad\\
I_q(BY)\le\\
\{E[h^4/q]\}^{1/3}\{E[B(|Y|^2)^2]\}^{2/3}.
\end{gathered}
\]

Therefore, for every real \(Y\),

\[
 I_q(BY)^3\le E[h^4/q]\,\|B\|^4(E|Y|^4)^2.
 \tag{Q87.2.1}\label{eq:Q87:2.1}
\]

For \(B=\allowbreak T\), \(h=\allowbreak q\), and
\(Eq^3=\allowbreak 1+3d-Eg^3\le\allowbreak 1+4d\). Together with (Q87.1.6), this proves
the first inequality of (Q87.0.3), since \(4+128/19=\allowbreak 204/19\).

For \(B=\allowbreak T^2\), the actual row is \(\bar f_2=\allowbreak q+r\). Expand the same row:

\[
 E[\bar f_2^4/q]
 =E[q^3+4q^2r]+E[(6q+4r+r^2/q)r^2].
\]

The first term is at most \(1+2d\): indeed it equals
\(1-d-Eg^3+4Eg^2r\), and (Q87.1.3) gives
\(4Eg^2r\le\allowbreak 2E[g^2pq]\le\allowbreak 2d\), while \(-Eg^3\le\allowbreak  d\).
For the second, (Q87.1.3) gives

\[
 6q+4r+r^2/q\le q(6+2p+p^2/4)\le12.
\]

The last polynomial is \(12-2p-3p^2/2-p^3/4\) for \(0\le\allowbreak  p\le\allowbreak 2\).
Consequently

\[
 E[\bar f_2^4/q]\le1+2d+12\delta_2\le1+55d/8.
 \tag{Q87.2.2}\label{eq:Q87:2.2}
\]

Now \(\|T^2\|^4\le\allowbreak (1+d/2)^8\le\allowbreak  e^{4d}\); (Q87.2.1) yields
the second inequality of (Q87.0.3). 

One further exact consequence will be used later:

\[
 I_q(T^3Y)^3\le(1-23d/36)(E|Y|^4)^2.
 \tag{Q87.2.3}\label{eq:Q87:2.3}
\]

This follows from (Q87.2.1), \(\|T^3\|\le\allowbreak 1\), and the actual
inverse-fourth row. 

\paragraph{Q87.3 Every unit inner connector with a singleton arm at least three}\label{part:Q87.3}

Use the actual reversed terminal \(F_b^-=\allowbreak p\bar f_b\). For \(b=\allowbreak 2\),

\[
 0<TF_2^-\le q-\gamma-s,\qquad \gamma=d-\delta_2/2.
 \tag{Q87.3.1}\label{eq:Q87:3.1}
\]

This follows from \(p\bar f_2\le\allowbreak (1+r/2)^2\), \(Tr=\allowbreak -d-s\),
and \(T(r^2)\le\allowbreak 2\delta_2\). For \(b\ge\allowbreak 3\), put

\[
\begin{gathered}
\bar R=KG_{b-1}(-K)g,\quad m=1-\beta_b/2,\quad k=1/m,\quad\\
t_0=\beta_{b+1},\quad e=d-\beta_{b+2},\quad V=E\bar R^2,
\end{gathered}
\]
\[
 \gamma=kt_0-k^2V/2.
\]

The exact positive comparison \eqref{eq:ordinary-unit-inner:terminal-long} is

\[
\begin{gathered}
0<TF_b^-\le m^2[q-\gamma-kK\bar R],\quad\\
1\le k\le6/5,\quad \gamma-ke\ge d-68\delta_2/25.
\end{gathered}\tag{Q87.3.2}\label{eq:Q87:3.2}
\]

Here \(E\bar R=\allowbreak -t_0\), \(E(g\bar R)=\allowbreak e\),
\(V=\allowbreak \delta_2+\delta_3\) for \(b=\allowbreak 3\), and \(V\le\allowbreak  \delta_2\) for \(b\ge\allowbreak 4\).
Both upper factors are nonnegative by their positive-kernel derivation.

Write \(R_a=\allowbreak KG_{a-1}(K)g\). The root envelope
\(f_a\le\allowbreak  p+R_a-\beta_a\) and AM–GM give

\[
\begin{gathered}
f_aT(p\bar f_b)\le\\
m^2[1-(\beta_a+\gamma)/2+Z/2]^2,\quad\\
Z=K[G_{a-1}(K)-kKG_{b-1}(-K)]g,
\end{gathered}\tag{Q87.3.3}\label{eq:Q87:3.3}
\]

with \(EZ=\allowbreak -\beta_{a+1}+ke\). For \(b=\allowbreak 2\), set \(m=\allowbreak k=\allowbreak 1\),
\(Z=\allowbreak K[G_{a-1}(K)-K]g\), \(EZ=\allowbreak -\beta_{a+1}\).
Every square base is half a sum of two nonnegative upper factors.
Thus raising to the eighth power and using \(1+x\le\allowbreak  e^x\) is legal.
Equation (Q87.1.5) at \(u=\allowbreak 4\) applies to this same fluctuation.

\paragraph{Q87.3.1 Terminal lengths at least three, root arms three and four}\label{part:Q87.3.1}

The exact correlated estimate \eqref{eq:ordinary-unit-inner:correlated-frequency} gives

\[
 \|[G_{a-1}(K)-kKG_{b-1}(-K)]g\|^2
 \le d-17\delta_2/20\qquad(a,b\ge3).
\]

Consequently

\[
\begin{gathered}
\log EY^4\le\\
-4(\beta_a+\beta_{a+1}+\beta_b+\gamma-ke)\\
+(17/4)d-(289/80)\delta_2.
\end{gathered}\tag{Q87.3.4}\label{eq:Q87:3.4}
\]

For \(a=\allowbreak 3\), \(\beta_3+\beta_4=\allowbreak 2d\); for \(a=\allowbreak 4\),
\(\beta_4+\beta_5=\allowbreak 2d-\delta_2\). Using \(\beta_b\ge\allowbreak \tau\) gives,
respectively,

\[
 \log EY^4\le-47d/4+(4507/400)\delta_2,\qquad
 \log EY^4\le-47d/4+(6107/400)\delta_2.
\]

The second is the weaker one, and its upper bound is
\(-71009d/12800<\allowbreak -11d/2\), because
\(150400-79391=\allowbreak 71009>\allowbreak 70400\). This proves (Q87.0.1) in both rows.

\paragraph{Q87.3.2 A stronger finite frequency bound for every root arm at least five}\label{part:Q87.3.2}

For \(m_0\ge\allowbreak 4,\allowbreak  n\ge\allowbreak 2,\allowbreak  1\le\allowbreak  k\le\allowbreak 6/5\), one has

\[
\begin{gathered}
|G_{m_0}(i\sqrt x)-ki\sqrt xG_n(-i\sqrt x)|^2\\
\le1-3x/2,\qquad0\le x\le\theta.
\end{gathered}\tag{Q87.3.5}\label{eq:Q87:3.5}
\]

To prove it write \(G_{m_0}=\allowbreak \alpha_0+i\sqrt x\beta_0\),
\(G_n(-i\sqrt x)=\allowbreak \gamma_0-i\sqrt x\omega_0\).
The first two coordinates lie in \([1-x,\allowbreak  1-x+x^2]\), the last
two in \([1-x,\allowbreak  1]\). The real part of the difference is positive,
at least \(1-(1+k)x\), and at most
\(1-(1+k)x+(1+k)x^2\). Its imaginary coefficient has magnitude
at most \(k-1+x\). Squaring gives \(1+xP(k,\allowbreak  x)\), where

\[
 P=k^2-4k-1+(k^2+6k+1)x
 -(2k^2+4k+1)x^2+(k+1)^2x^3.
\]

Its derivative in \(x\) is at least
\(8-(434/25)\theta=\allowbreak 379/400>\allowbreak 0\). Its second derivative in \(k\)
is \(2(1+x-2x^2+x^3)>\allowbreak 0\). The two endpoint maxima at \(x=\allowbreak \theta\)
are

\[
 P(1,\theta)=-13411/8192<-3/2,\qquad
 P(6/5,\theta)=-1271219/819200<-3/2.
\]

These prove the entire continuous rectangle and all integer lengths.

For \(a\ge\allowbreak 5\), retain also the actual term
\(\beta_{a-1}p\ge\allowbreak \tau p\) in (Q87.1.4):

\[
 f_a\le(1-\tau)p+R_a-\beta_a.
\]

In (Q87.3.3), the new deterministic loss is \(\beta_a+\gamma+\tau\)
and the fluctuation is \(Z'=\allowbreak K(V_0-\tau1)\), where
\(V_0=\allowbreak [G_{a-1}(K)-kKG_{b-1}(-K)]g\).
Its mean is unchanged. Moreover

\[
 EV_0=-\beta_a+kt_0\ge-\delta_2,\qquad
 \|V_0-\tau1\|^2\le d-3\delta_2/2+d^2-\delta_2^2.
\]

Using (Q87.1.5), (Q87.3.2), and \(\beta_a,\allowbreak  \beta_{a+1},\allowbreak  \beta_b\ge\allowbreak \tau\),

\[
\begin{gathered}
\log EY^4\le-63d/4+(4101/200)\delta_2+(17/4)(d^2-\delta_2^2)\\
\le-115261d/19200<-6d.
\end{gathered}\tag{Q87.3.6}\label{eq:Q87:3.6}
\]

The last step only uses \(\delta_2\le\allowbreak \theta d\), \(d^2\le\allowbreak  d/3\),
and discards the negative \(\delta_2^2\). Thus every \(a\ge\allowbreak 5,\allowbreak  b\ge\allowbreak 3\)
is covered without a remaining score interval.

\paragraph{Q87.3.3 Terminal two, every root arm}\label{part:Q87.3.3}

For \(b=\allowbreak 2\), (Q87.3.3) and (Q87.1.5) give

\[
\begin{gathered}
\log EY^4\le-4(\beta_a+\beta_{a+1}+\gamma)\\
+(17/4)\|[G_{a-1}(K)-K]g\|^2.
\end{gathered}\tag{Q87.3.7}\label{eq:Q87:3.7}
\]

At \(a=\allowbreak 3\), the last polynomial is 1; this yields
\(-31d/4+2\delta_2\le\allowbreak -111d/16\).
At \(a=\allowbreak 4\), its squared modulus is \((1-x)^2\), and (Q87.3.7)
is at most \(-31d/4-(5/2)\delta_2+(17/4)\delta_3\le\allowbreak -31d/4\).

At \(a=\allowbreak 5\), the exact squared modulus is
\(1-2x+x^2+x^3\), while \(\beta_5+\beta_6=\allowbreak 2\tau\).
The upper coefficient is

\[
 -31/4+(3/2)\theta+(17/4)\theta^2(1+\theta)<-91/16.
\]

For an explicit comparison, the last positive term is
\(129285/131072<\allowbreak 1\); the whole expression is less than
\(-31/4+39/64+1=\allowbreak -393/64<\allowbreak -91/16\).

At \(a=\allowbreak 6\), the exact squared modulus is
\(1-2x+3x^2-x^3+x^4\), and
\(\beta_6+\beta_7=\allowbreak 2\tau+\delta_3\). Hence the bound is at most
\(-31/4+(3/2)\theta+(35/4)\theta^2=\allowbreak -23333/4096<\allowbreak -91/16\).

For every \(a\ge\allowbreak 7\), the same upper squared polynomial
\(1-2x+3x^2-x^3+x^4\) follows from the finite coordinate bounds.
Also \(\beta_a,\allowbreak  \beta_{a+1}\ge\allowbreak \tau+\kappa \delta_3\), by the finite
alternating tails. Substitution in (Q87.3.7) gives

\[
 \log EY^4\le-31d/4+(3/2)\delta_2+8\delta_3
 -(17/4)\delta_4+(17/4)\delta_5
 \le-745d/128<-91d/16.
\]

All terminal-two cases therefore satisfy (Q87.0.1). This completes every
\(a\ge\allowbreak 3,\allowbreak  b\ge\allowbreak 2,\allowbreak  L=\allowbreak 1\), with the single orientation in the local bound.

\paragraph{Q87.4 The nonunit-inner scalar table}\label{part:Q87.4}

The remaining proof uses the scalar estimates in Section~\ref{proof:Q5-C075}, the scalar-input and row-table estimates.
Set \(X=\allowbreak T^L(qf_b)\), \(L\ge\allowbreak 2\), and \(W_2=\allowbreak p-r\). The
positive envelope has the form

\[
 0<X\le C[Bq+Br+Z-a_0],\qquad
 W_2X\le CB[1+(Z-a_0)/(2B)]^2.
 \tag{Q87.4.1}\label{eq:Q87:4.1}
\]

Every \(B\) below is positive. If \(Z=\allowbreak K^2V_0\), the
paid expression is

\[
 S_8=8\log C+8\log B-8a_0/B+8M_0/B+6D/B^2.
 \tag{Q87.4.2}\label{eq:Q87:4.2}
\]

Here \(EZ\le\allowbreak  M_0\), \(\|V_0\|^2\le\allowbreak  D\). For \(Z=\allowbreak KV_0\),
replace its last coefficient 6 by \(8c_*\),
\(c_*=\allowbreak (1+\theta)^2=\allowbreak 2025/1024\). The exact rows are:


\begingroup\small
\setlength{\tabcolsep}{3pt}
\setlength{\LTleft}{0pt}\setlength{\LTright}{\fill}
\begin{longtable}{@{}p{0.2200\linewidth}p{0.7000\linewidth}@{}}
\caption{Ordinary fourth moments for genuine singleton-root inputs: parameter values.}\label{tab:parameters:q5-c087:2}\\
\toprule
Entry & \(C,\allowbreak  B,\allowbreak  a_0\) \\
\midrule\endfirsthead
\toprule
Entry & \(C,\allowbreak  B,\allowbreak  a_0\) \\
\midrule\endhead
\bottomrule\endfoot
\(b=\allowbreak 2,\allowbreak  L=\allowbreak 2\) & \(1,\allowbreak  1-\gamma,\allowbreak  0\); Row vector and \(D\): two-row \(V_0=\allowbreak \gamma1+r\), \(D=\allowbreak \delta_2\); Accepted upper for \(S_8\): \(-d/8\). \\
\addlinespace[2pt]
\(b=\allowbreak 2,\allowbreak  L=\allowbreak 3\) & \(1,\allowbreak  1-\gamma,\allowbreak  \tau\); Row vector and \(D\): two-row \(V_0=\allowbreak -(g+s)\), \(D=\allowbreak d-2\delta_2+\delta_3\); Accepted upper for \(S_8\): \(-3d/4\). \\
\addlinespace[2pt]
\(b=\allowbreak 2,\allowbreak  L=\allowbreak 4\) & \(1,\allowbreak  1-\gamma-\tau,\allowbreak  d(1-\gamma)\); Row vector and \(D\): two-row \(V_0=\allowbreak \tau1-(1-\gamma)g+r+t\); Accepted upper for \(S_8\): \(-d\). \\
\addlinespace[2pt]
\(b=\allowbreak 2,\allowbreak  L=\allowbreak 5\) & \(1,\allowbreak  1-2\tau,\allowbreak  \tau\); Row vector and \(D\): one-row \([\tau+zP_3(z)]g\), \(D/d=\allowbreak \theta-2\theta\tau+\tau^2\); Accepted upper for \(S_8\): \(-2d/3\). \\
\addlinespace[2pt]
\(b=\allowbreak 2,\allowbreak  L\ge\allowbreak 6\) & \(1,\allowbreak  1-2\tau,\allowbreak  \tau\); Row vector and \(D\): one-row \([\tau+zP_{L-2}(z)]g\), \(D/d=\allowbreak \theta-7\theta^2/15-2\theta\kappa\tau+\tau^2\); Accepted upper for \(S_8\): \(-4d/5\). \\
\addlinespace[2pt]
\(b\ge\allowbreak 3,\allowbreak  L=\allowbreak 2\) & \(m^2,\allowbreak  1-\gamma,\allowbreak  ke\); Row vector and \(D\): two-row \(V_0=\allowbreak \gamma1+kR\), \(D=\allowbreak k^2V\); Accepted upper for \(S_8\): \(-d/8\). \\
\addlinespace[2pt]
\(b\ge\allowbreak 3,\allowbreak  L=\allowbreak 3\) & \(m^2,\allowbreak  1-B_*,\allowbreak  d-k\rho\); Row vector and \(D\): one-row \([\delta+zP_1(z)]g\), \(D/d=\allowbreak 351/1280+\delta^2\), \(\delta=\allowbreak ke\); Accepted upper for \(S_8\): \(-d/10\). \\
\addlinespace[2pt]
\(b\ge\allowbreak 3,\allowbreak  L=\allowbreak 4\) & \(m^2,\allowbreak  1-41d/80,\allowbreak  \kappa d\); Row vector and \(D\): one-row \([\delta+zP_2(z)]g\), \(D/d=\allowbreak \theta(71/63)+\delta^2\), \(\delta=\allowbreak 41d/80\); Accepted upper for \(S_8\): \(-8d/25\). \\
\addlinespace[2pt]
\(b\ge\allowbreak 3,\allowbreak  L\ge\allowbreak 5\) & \(1-d/2,\allowbreak  1-d/2,\allowbreak  d/2\); Row vector and \(D\): one-row \([\delta+zP_{L-2}(z)]g\), \(D/d=\allowbreak \theta+d^2/4\), \(\delta=\allowbreak d/2\); Accepted upper for \(S_8\): \(-2d/5\). \\
\end{longtable}
\endgroup


For \(b=\allowbreak 2\), \(\gamma=\allowbreak d-\delta_2/2\), \(z=\allowbreak -K\),
\(P_n=\allowbreak G_n+z^{n+1}\). In its \(L=\allowbreak 4\) row,
\(D=\allowbreak (1-\gamma)^2d+\delta_2-2\delta_3+\delta_4-\tau^2\).
For \(b\ge\allowbreak 3\), use the actual positive terminal remainder
\(R=\allowbreak KG_{b-1}(K)g\) and the same even-moment quantities
\(m,\allowbreak  k,\allowbreak  t_0,\allowbreak  e,\allowbreak  V,\allowbreak  \gamma\) as in (Q87.3.2), with
\(\rho=\allowbreak ErR\), \(B_*=\allowbreak \gamma+ke\), and
\(P_n=\allowbreak G_n+kz^{n+1}G_{b-1}(-z)\).

The same-parameter bounds are

\[
\begin{gathered}
3d/10\le\gamma\le6d/5,\quad\\
5d/8\le B_*\le3d/2,\quad\\
d-k\rho\ge41d/80,\quad 0\le ke\le39d/80,
\end{gathered}
\]
\[
\begin{gathered}
\kappa \delta_2\le\rho\le \delta_2,\quad\\
m^2\le1-209d/384,\quad 1-\gamma\ge17/25.
\end{gathered}\tag{Q87.4.3}\label{eq:Q87:4.3}
\]

For \(b=\allowbreak 2\), the bounds needed here are
\(1-\gamma\ge\allowbreak 13/18\), \(1-\gamma-\tau\ge\allowbreak 1/2\),
\(1-2\tau\ge\allowbreak 5/9\). In the three \(b\ge\allowbreak 3,\allowbreak  L\ge\allowbreak 3\) rows,
\(M_0=\allowbreak (\theta-\delta)d\). In the two \(L=\allowbreak 2\) rows,
\(M_0=\allowbreak \max(k\rho-\gamma d,\allowbreak  0)\), with \(k=\allowbreak 1,\allowbreak  \rho=\allowbreak \delta_2\) for
\(b=\allowbreak 2\); that latter mean is already nonnegative.
For \(b=\allowbreak 2,\allowbreak  L=\allowbreak 3\), \(M_0=\allowbreak 0\); for \(L=\allowbreak 4\),
\(M_0=\allowbreak (\theta\kappa-\tau)d\).

All row-vector ranges, raw second moments, and means belong to these
same actual fluctuations. The \(b\ge\allowbreak 3,\allowbreak  L=\allowbreak 2\) scalar includes the terminal end-sign term from Section~\ref{proof:Q5-C071}/Section~\ref{proof:Q5-C075}. Every row has \(S_8\le\allowbreak -d/10\).

For later covariance arguments, write
\(P_n=\allowbreak A_P+zB_P\), \(z=\allowbreak i\sqrt x\). The finite coefficient
lemma says

\[
 1-x\le A_P,B_P\le1\quad(n\ge3),
\]
\[
\begin{gathered}
n=1: A_P\ge1-6x/5, 0\le B_P\le6x/5;\qquad\\
n=2: A_P\ge1-x, 1-6x/5\le B_P\le1.
\end{gathered}\tag{Q87.4.4}\label{eq:Q87:4.4}
\]

These include every finite tail. In particular \(\Re P_n\ge\allowbreak 1/2\)
in all required rows.

\paragraph{Q87.5 Every nonunit inner connector with a singleton arm at least five}\label{part:Q87.5}

Let \(a\ge\allowbreak 5\), \(m_0=\allowbreak a-2\ge\allowbreak 3\), \(U_0=\allowbreak K^2G_{m_0}(-K)g\),
and \(\beta=\allowbreak \beta_a\). Apply positive \(A\) to (Q87.1.4) for
\(\bar f_{a-1}\). Its last two nonconstant convolution terms are the
actual messages \(\beta_{a-1}p\) and \(\beta_{a-2}W_2\).
Retaining both, and using their actual lower coefficients \(\tau\),
gives

\[
 0<W_a\le cp-(1-\tau)r+U_0-\beta,\qquad c=1-2\tau.
 \tag{Q87.5.1}\label{eq:Q87:5.1}
\]

This upper side is nonnegative because it equals the actual path
plus omitted nonnegative path terms and nonnegative coefficient
remainders. 

Set \(H=\allowbreak G_{m_0}(-K)g-\tau1\). Scaled AM–GM with (Q87.4.1) gives

\[
 W_aX\le CBc
 [1+(Z+(B/c)K^2H-a_0-(B/c)\beta)/(2B)]^2.
 \tag{Q87.5.2}\label{eq:Q87:5.2}
\]

Indeed \(K^2H=\allowbreak U_0-\tau r\); the coefficient of \(r\) after
dividing the two factors by \(c\) and \(CB\) is exactly \(-\tau/c\).
Both scaled factors are nonnegative.

The exact Section~\ref{proof:Q5-C075} centered-covariance argument gives

\[
 EK^2H\le0,\qquad \|H\|^2\le d-\tau^2,
 \tag{Q87.5.3}\label{eq:Q87:5.3}
\]

and the nonpositive original inner products
\(\langle V_0,\allowbreak  H\rangle\le\allowbreak 0\) in every two-row case, or
\(\langle V_0,\allowbreak  KH\rangle\le\allowbreak 0\) in every one-row case.
These statements were proved in Section~\ref{proof:Q5-C075}, the centered-covariance estimates for this same
\(H\), all \(m_0\ge\allowbreak 2\), and every row of part Q87.4. They are independent of \(c\) and therefore apply in (Q87.5.2). In particular the Section~\ref{proof:Q5-C075} proof retains
\(\|H\|^2\le\allowbreak  d-\tau^2\), rather than its later weaker \(\le\allowbreak  d\).

For clarity, the mean uses
\(\delta_2-\delta_3-\tau d\ge\allowbreak (\kappa-d)(\delta_2-d^2)\ge\allowbreak 0\).
The latter follows by applying the original spectral bound to
\(r+d1\). The covariance at \(L=\allowbreak 2\) is the actual covariance of
the two finite arm remainders; the other one-row covariances are
the finite coefficient inequalities in Section~\ref{proof:Q5-C075}. 

Put \(\alpha=\allowbreak 17\theta/4=\allowbreak 221/128\). Apply (Q87.1.5) at \(u=\allowbreak 4/B\)
to (Q87.5.2). In a two-row case its variance coefficient is at most
\(\alpha\|V_0+(B/c)H\|^2/B^2\). In a one-row case it is at
most \((17/4)\|V_0+(B/c)KH\|^2/B^2\).
The covariance signs, (Q87.5.3), and \(\|K\|^2\le\allowbreak \theta\) imply

\[
 \log E[(W_aX)^4]\le S_8/2+F(d,\tau),
\]
\[
\begin{gathered}
F(d,\tau)=4\log(1-2\tau)-\frac{4\tau}{1-2\tau}\\
+\alpha\frac{d-\tau^2}{(1-2\tau)^2}.
\end{gathered}\tag{Q87.5.4}\label{eq:Q87:5.4}
\]

Here \(\alpha<\allowbreak 3\) and \(17/4<\allowbreak 4c_*\) pay the old vector costs.
The old mean upper stays valid because \(EK^2H\le\allowbreak 0\).

The following pays (Q87.5.4) on the entire original domain:

\[
 F(d,\tau)\le-691d/128.
 \tag{Q87.5.5}\label{eq:Q87:5.5}
\]

For fixed \(d\), its \(\tau\)-derivative, multiplied by
\((1-2\tau)^3\), is

\[
 -12+(40-2\alpha)\tau-32\tau^2+4\alpha d
 \le-227/144<0
\]

using \(\tau\le\allowbreak 2/9,\allowbreak  d\le\allowbreak 1/3\). Thus it suffices to put
\(\tau=\allowbreak \kappa d\). Let \(h=\allowbreak 12\kappa-\alpha=\allowbreak 691/128\).
Using \(\log(1-2\tau)\le\allowbreak -2\tau-2\tau^2\), the assertion
\(F+hd\le\allowbreak 0\), after multiplication by \((1-2\tau)^2\), is
bounded above by

\[
 \tau^2[32-\alpha-4h/\kappa+4hd-32\kappa^2d^2].
\]

The bracket increases on \([0,\allowbreak  1/3]\), since its derivative is at
least \(1351/96>\allowbreak 0\). Its endpoint is

\[
 3875/128-691/19+691/96-361/288=-3295/21888<0.
\]

This proves (Q87.5.5).
Since \(S_8\le\allowbreak -d/10\), (Q87.5.4) now gives precisely (Q87.0.2).

\paragraph{Q87.6 Every nonunit inner connector with a singleton arm four}\label{part:Q87.6}

Here the exact coefficient retained from (Q87.1.4) is \(\beta_3=\allowbreak d\),
not merely its lower bound \(\tau\). The root identity is

\[
 W_4=(1-d)p-r+(s-t)-d.
\]

Put \(c=\allowbreak 1-d\), \(H=\allowbreak g-r-d1\). Formula (Q87.5.2) holds with
\(\beta=\allowbreak d\), and

\[
 EH=0,\quad\|H\|^2=d+\delta_2-d^2,\quad EK^2H=-\delta_2+d^2\le0.
 \tag{Q87.6.1}\label{eq:Q87:6.1}
\]

The required two-row covariances remain nonpositive. At \(L=\allowbreak 2\),
the covariance proof in Section~\ref{proof:Q5-C075}, the two-row covariance calculation has no centering constant
inside its covariance, and \(EH=\allowbreak 0\). At \(b=\allowbreak 2,\allowbreak  L=\allowbreak 3,\allowbreak  4\), \(EV_0=\allowbreak 0\),
so replacing \(\tau1\) by \(d1\) leaves the covariance
unchanged.

Every one-row covariance is also nonpositive
with this new centering. At \(z=\allowbreak -K=\allowbreak i\sqrt x\), \(G_2(z)=\allowbreak 1+z\),
and the spectral product controlling its negative is

\[
 \delta d+x[A_P-\delta+(x-d)B_P].
 \tag{Q87.6.2}\label{eq:Q87:6.2}
\]

For \(n\ge\allowbreak 3\), and also \(n=\allowbreak 2\), the square bracket is at least
\(\kappa-2/9-1/3=\allowbreak 11/288>\allowbreak 0\), by (Q87.4.4).
For \(n=\allowbreak 1\), it is at least
\(41/80-13/80-13/80=\allowbreak 15/80>\allowbreak 0\).
Thus the same comparison to \(S_8/2\) gives the increment

\[
 F_4=4\log(1-d)-4d/(1-d)
 +\alpha[(1+\theta)d-d^2]/(1-d)^2.
\]

Let \(h_4=\allowbreak 8-\alpha(1+\theta)=\allowbreak 22823/4096\).
Using \(\log(1-d)\le\allowbreak -d-d^2/2\), multiplication by \((1-d)^2\)
bounds \(F_4+h_4d\) by

\[
 d^2[10-\alpha-2h_4+h_4d-2d^2]
 \le d^2[10-\alpha-5h_4/3]
 =-12451d^2/12288\le0.
\]

Hence \(\log E[(W_4X)^4]\le\allowbreak -(h_4+1/20)d<\allowbreak -11d/2\).
This covers every terminal length and every \(L\ge\allowbreak 2\).

\paragraph{Q87.7 Singleton arm three and every inner length at least three}\label{part:Q87.7}

The exact root is \(W_3=\allowbreak W_2+s-d\). Combining with (Q87.4.1),

\[
 W_3X\le CB[1+(Z+Bs-a_0-Bd)/(2B)]^2.
 \tag{Q87.7.1}\label{eq:Q87:7.1}
\]

The new fluctuation has the same mean as \(Z\); its deterministic
loss adds \(-4d\) to the ordinary-fourth exponent.
For one-row cases its row vector becomes
\([\delta+z(P_n-B)]g\). At every original frequency,

\[
 |\delta+z(P_n-B)|^2-|\delta+zP_n|^2
 =Bx(B-2\Re P_n)\le0.
 \tag{Q87.7.2}\label{eq:Q87:7.2}
\]

Put \(\Delta=\allowbreak 4c_*-17/4=\allowbreak 937/256\). Except for the separately
treated \(b\ge\allowbreak 3,\allowbreak  L=\allowbreak 3\) row, (Q87.1.5), (Q87.7.2), and the paid table yield

\[
 \log E[(W_3X)^4]\le S_8/2-4d-\Delta D/B^2
 \tag{Q87.7.3}\label{eq:Q87:7.3}
\]

in every remaining one-row case.

For \(b=\allowbreak 2,\allowbreak  L=\allowbreak 5\), \(D/B^2\ge\allowbreak \theta d\). For \(b=\allowbreak 2,\allowbreak  L\ge\allowbreak 6\),
it is at least \((\theta-7\theta^2/15)d\). To verify both entire
domains, put \(A(t)=\allowbreak A_0-2\theta\sigma t+t^2\). The derivative
of \(A(t)/(1-2t)^2\) has positive numerator
\(4A_0-2\theta\sigma+2t(1-2\theta\sigma)\), for
\((A_0,\allowbreak  \sigma)=\allowbreak (\theta,\allowbreak  1)\) or
\((\theta-7\theta^2/15,\allowbreak  \kappa)\). Its minimum is therefore at zero.
For \(b\ge\allowbreak 3,\allowbreak  L=\allowbreak 4\), \(D/B^2\ge\allowbreak \theta(71/63)d>\allowbreak \theta d\);
for \(b\ge\allowbreak 3,\allowbreak  L\ge\allowbreak 5\), \(D/B^2\ge\allowbreak \theta d\).

The exact lower constants used in (Q87.7.3) are

\[
\begin{gathered}
\Delta\theta=12181/8192>67/50,\qquad\\
\Delta(\theta-7\theta^2/15)=4738409/3932160>11/10.
\end{gathered}
\]

With the corresponding \(S_8/2\) values, each coefficient is
strictly more negative than \(-11d/2\). In particular the weakest
long terminal-two row has loss at least \(4+2/5+11/10=\allowbreak 11/2\).

For \(b=\allowbreak 2,\allowbreak  L=\allowbreak 3\), the new two-row vector is \(-(\gamma g+s)\).
Since \(\gamma d\le\allowbreak  \delta_2\), its square norm is at most \(\delta_3\).
Its mean is zero, and direct use of (Q87.1.5) gives

\[
\begin{gathered}
\log E[(W_3X)^4]\\
\le4\log(1-\gamma)-4\tau/(1-\gamma)-4d\\
+\alpha \delta_3/(1-\gamma)^2\\
\le[-153/16+17901/32768]d<-9d.
\end{gathered}\tag{Q87.7.4}\label{eq:Q87:7.4}
\]

For \(b=\allowbreak 2,\allowbreak  L=\allowbreak 4\), the new vector is \(\tau1-\tau g+r+t\), whose
norm square is
\(\delta_2-2\delta_3+\delta_4-\tau^2+\tau^2d\le\allowbreak  \delta_2+\tau^2d\).
Using the same table parameters gives

\[
\begin{gathered}
\log E[(W_3X)^4]\le\\
-(89/16)d-4d-4d+8\theta\kappa d\\
+4\alpha(\theta+4/81)d<-8d.
\end{gathered}\tag{Q87.7.5}\label{eq:Q87:7.5}
\]

For instance \(8\theta\kappa<\allowbreak 2\), while
\(4\alpha(\theta+4/81)=\allowbreak 2873/1024+221/648<\allowbreak 13/4\), so the last
strict comparison follows on the entire domain.

Finally let \(b\ge\allowbreak 3,\allowbreak  L=\allowbreak 3\), \(t_*=\allowbreak B_*\in[0,\allowbreak  1/2]\), \(B=\allowbreak 1-t_*\).
The new one-row vector is \([\delta+z(P_1-B)]g\). Writing
\(G_{b-1}(-z)=\allowbreak A-zD_0\), its squared modulus is

\[
\begin{gathered}
(\delta-kx^2D_0)^2+x(t_*-kxA)^2\\
\le\delta^2+\theta t_*^2+c_0,\qquad\\
c_0=(36/25)\theta^3(1+\theta)<7/50.
\end{gathered}
\]

The last rational bound follows from
\(\theta^3<\allowbreak 27/400\), \(1+\theta<\allowbreak 10/7\), and
\((36/25)(27/400)(10/7)=\allowbreak 243/1750<\allowbreak 7/50\).
The old row has \(D/d=\allowbreak 351/1280+\delta^2\). Therefore its saved
cost, divided by \(d\), is at least

\[
 \frac{3/2-\alpha t_*^2}{(1-t_*)^2}\ge3/2.
\]

Indeed
\[
4c_*(351/1280)-(17/4)c_0
>(79/10)(27/100)-(17/4)(7/50)=769/500>3/2,
\]
and the displayed quotient increases on \([0,\allowbreak  1/2]\), since its
derivative has numerator \(2(3/2-\alpha t_*)>\allowbreak 0\).
Consequently its fourth exponent is at most
\(-d/20-4d-3d/2=\allowbreak -111d/20<\allowbreak -11d/2\).
This finishes every \(L\ge\allowbreak 3\) at \(a=\allowbreak 3\).

\paragraph{Q87.8 Singleton arm three, two original inner edges, and a real turn}\label{part:Q87.8}

For the remaining short-outer row use

\[
 Y=f_3TA(p\bar f_b).
 \tag{Q87.8.1}\label{eq:Q87:8.1}
\]

The two inner edges meet at their original midpoint and are both
directed away from that midpoint. Thus at the singleton root the
last-edge operator is \(T\), and at the midpoint the terminal-edge
operator is \(A\). The root arm is the outward three-path \(f_3=\allowbreak p+r+s-d\).

The positive terminal square used in \hyperref[part:Q87.3]{part Q87.3} gives

\[
 A(p\bar f_b)\le m^2[p-\gamma+kK\bar R],
\]
\[
 0<TA(p\bar f_b)\le m^2[Bq-r-kK^2\bar R-ke],\qquad B=1-\gamma.
 \tag{Q87.8.2}\label{eq:Q87:8.2}
\]

For \(b=\allowbreak 2\), use \(m=\allowbreak k=\allowbreak 1,\allowbreak  \bar R=\allowbreak r,\allowbreak  e=\allowbreak 0\). Set
\(V_0=\allowbreak \gamma1+k\bar R\), \(\rho=\allowbreak Er\bar R\).
Scaled AM–GM yields

\[
 Y\le m^2B[1-d/2-ke/(2B)+Z/2]^2,\qquad
 Z=K^2[g-V_0/B].
 \tag{Q87.8.3}\label{eq:Q87:8.3}
\]

Its exact mean and row-vector norm satisfy

\[
\begin{gathered}
EZ=(\gamma d-k\rho)/B,\qquad\\
\|g-V_0/B\|^2\\
\le d-2ke/B+k^2V/B^2.
\end{gathered}\tag{Q87.8.4}\label{eq:Q87:8.4}
\]

Here \(E(g\bar R)=\allowbreak e\), while \(\|V_0\|^2\le\allowbreak  k^2V\) follows
from the identity \eqref{eq:Q71:8.2}. The sign of the mean
is retained; (Q87.1.5) allows either sign. Thus

\[
\begin{gathered}
\log EY^4\le4\log m^2+4\log B-4d+\alpha d\\
+[4\gamma d-4k\rho-(4+2\alpha)ke]/B\\
+\alpha k^2V/B^2.
\end{gathered}\tag{Q87.8.5}\label{eq:Q87:8.5}
\]

\paragraph{Q87.8.1 Terminal two}\label{part:Q87.8.1}

Here \(\rho=\allowbreak \delta_2\), \(B\ge\allowbreak 13/18\), and the exact norm is

\[
 \|V_0\|^2=\delta_2-d^2+\delta_2^2/4.
\]

Put \(c_0=\allowbreak \alpha/(13/18)^2=\allowbreak 1377/416\).
Use \(\delta_2-\gamma d\ge\allowbreak 0\), \(B\le\allowbreak 1\), and \(\log B\le\allowbreak -\gamma\)
in (Q87.8.3)–(Q87.8.5). After collecting the same moments,

\[
 \log EY^4\le-(8-\alpha)d+(c_0-2)\delta_2+(4-c_0)d^2
 -2d\delta_2+(c_0/4)\delta_2^2.
\]

The last two terms sum to a nonpositive number because
\((c_0/4)\delta_2\le\allowbreak (c_0\theta/4)d<\allowbreak 2d\).
Since \(c_0>\allowbreak 2\) and \(4-c_0>\allowbreak 0\), the full bound is at most

\[
 [-8+\alpha+(c_0-2)\theta+(4-c_0)/3]d
 =-220097d/39936<-11d/2.
 \tag{Q87.8.6}\label{eq:Q87:8.6}
\]

The final numerator comparison is \(220097>\allowbreak 219648\).

\paragraph{Q87.8.2 Every terminal length at least four}\label{part:Q87.8.2}

Use the same \(\delta_2\) throughout, and put

\[
 C_0=1989/800,\quad N_0=13927/2048,\quad
 \gamma_0=d-43\delta_2/25>0.
\]

The finite alternating sums imply
\(\beta_b\ge\allowbreak \tau\), \(e,\allowbreak  \rho\ge\allowbreak \kappa \delta_2\),
\(V\le\allowbreak  \delta_2\), and \(\gamma\ge\allowbreak \gamma_0\).
Thus (Q87.8.5) is bounded by

\[
\begin{gathered}
-4\tau-4d+\alpha d+\\
\Phi(B),\quad\\
\Phi(B)=4\log B+4(1-B)d/B+C_0\delta_2/B^2-N_0\delta_2/B.
\end{gathered}
\]

On \(B\ge\allowbreak 17/25\), its derivative is positive. After multiplying
by \(B^3\), the derivative is
\(4B^2-4dB+N_0\delta_2B-2C_0\delta_2\).
The first two terms are at least \(1768/1875>\allowbreak 9/10\);
\(2C_0\delta_2<\allowbreak 5(7/50)=\allowbreak 7/10\); the omitted term is nonnegative.
Hence we may replace \(B\) by \(B_0=\allowbreak 1-\gamma_0\ge\allowbreak  B\).
Now \(2/3\le\allowbreak  B_0\le\allowbreak 1\).

For \(1\le\allowbreak  w\le\allowbreak 3/2\), convexity and \(N_0\ge\allowbreak 5C_0/2\) imply
\(C_0w^2-N_0w\le\allowbreak  C_0-N_0\).
Also \(4\gamma_0d/B_0\le\allowbreak 6d^2\). Using
\(\log B_0\le\allowbreak -\gamma_0\) gives

\[
\begin{gathered}
\log EY^4\le(-10+\alpha)d+\\
\relax [272/25+C_0-N_0]\delta_2\\
\le-9184899d/1638400<-28d/5.
\end{gathered}\tag{Q87.8.7}\label{eq:Q87:8.7}
\]

The final comparison is \(9184899>\allowbreak 9175040\).


\paragraph{Q87.8.3 Terminal three}\label{part:Q87.8.3}

Here \(\beta_b=\allowbreak d\), \(e=\allowbreak \rho=\allowbreak \delta_2\), \(V=\allowbreak \delta_2+\delta_3\).
Repeat the preceding argument with

\[
\begin{gathered}
C_0=\alpha(36/25)(1+\theta)<7/2,\quad\\
N_0=8+2\alpha>11,\quad\\
\gamma_0=d-81\delta_2/80.
\end{gathered}
\]

Monotonicity now follows immediately from
\(N_0B-2C_0>\allowbreak 11(17/25)-7>\allowbreak 0\) and \(B>\allowbreak d\).
The convex endpoint bound still holds since \(N_0\ge\allowbreak 5C_0/2\).
The final upper is

\[
 (-10+\alpha)d+[81/20+C_0-N_0]\delta_2
 \le(-10+\alpha)d<-11d/2.
\]

This completes (Q87.0.1) for (Q87.8.1). For \(M\ge\allowbreak 3\), the direction table instead uses \(W_3T^2(qf_b)\) and its eighth-moment estimate.
